\documentclass[10pt,a4paper]{amsart}
\usepackage[margin=19mm]{geometry}
\usepackage{amsmath,amssymb,mathtools}
\usepackage[scr=rsfs,cal=euler]{mathalfa}
\usepackage{xfrac}
\usepackage[unicode,hidelinks,bookmarks=false]{hyperref}
\newcommand{\ie}{i.e.\ }
\newcommand{\cf}{cf.\ }
\newcommand{\resp}{respectively\ }
\newcommand{\Subsection}{\subsection}
\providecommand{\nobreakdash}{\nobreak\hbox{-}}
\setlength{\emergencystretch}{2em}
\multlinegap0pt
\usepackage{bm}
\usepackage{bbm}
\usepackage{dsfont}
\usepackage{xspace}
\usepackage{cancel}
\usepackage{mathtools}
\usepackage{amsthm}
\makeatletter
\@ifundefined{theorem}{\newtheorem{theorem}{Theorem}}{}
\@ifundefined{lemma}{\newtheorem{lemma}[theorem]{Lemma}}{}
\makeatother
\newcommand{\Gpres}[2]{\bigl\langle #1\:|\:#2 \bigr\rangle}
\newcommand{\Mgen}[1]{\operatorname{Mon}\bigl\langle #1 \bigr\rangle}
\newcommand{\Z}{\mathbb{Z}}
\title[The word problem for two-generator one-relator inverse monoi]{The word problem for two-generator one-relator inverse monoids}
\author{Robert D. Gray, Catherine Reilly}
\date{Source: arXiv:2608.04650v1 (CC BY 4.0)}
\begin{document}
\maketitle

\noindent\emph{Source and licence.} The word problem for two-generator one-relator inverse monoids. Version arXiv:2608.04650v1 (CC BY 4.0), distributed under the Creative Commons Attribution 4.0 International licence, as recorded in the arXiv licence metadata for this exact version and shown on the abstract page https://arxiv.org/abs/2608.04650v1. Full rights notice: licenses/arxiv-2608.04650-CC-BY-4.0.txt.\par\smallskip

\noindent\textit{Editorial scope.} Counted unit: the Lemma whose source label is \texttt{la3} in the article (source0.tex lines 676--681, source label \texttt{la3}), delivered together with the article's own complete original proof of it (source0.tex lines 682--710, one \texttt{proof} environment of 29 lines). No mathematics is written, repaired or re-derived: statement and proof are the article's own text line for line. Required context retained uncounted: lem:APinfty (lines 416--420); lem:UsefulTExpSum (lines 445--450); thm:PrefixTech (lines 464--474); la2 (lines 497--503); thm:PrefixTech:Reduced (lines 618--640); lem:PrefGenSet:Reduced (lines 652--658). Every definition, display and label of those lines is retained unchanged. Omitted from this package and not redistributed: 1--415, 421--444, 451--463, 475--496, 504--617, 641--651, 659--675, 711--1426. Nothing inside a retained excerpt is omitted or altered: every retained body is delivered exactly as the article's lines are written, so no omission receipt of any kind accompanies this package. Citations: the retained excerpts carry no citation command at all, so no bibliography entry of the article is transcribed and no reference is redistributed. Reference adaptation: none was needed and none was made; every reference inside the delivered unit resolves inside it, so no unresolved reference or citation is produced. Printed slips kept literal: no printed word, symbol, hypothesis or number of the counted statement or of its proof was altered. Layout and class: the article's own class is \texttt{amsart} with packages including thmtools, thm-restate, pdfsync, enumerate, todonotes, mathabx, epsfig, tikz, tikz-cd, hyperref, clrscode, mathrsfs, amssymb, amsmath; the wrapper is the lane's pinned \texttt{amsart} editorial wrapper at 10pt on A4, which restates the theorem-like environments the retained text needs and adds the article's own macro definitions that the retained text needs (\texttt{\textbackslash Gpres}, \texttt{\textbackslash Mgen}, \texttt{\textbackslash Z}), with extra wrapper packages (bm, bbm, dsfont, xspace, cancel, mathtools). The delivered numbering is the unit's own. Assets: no retained span contains a figure environment, a graphics-inclusion command, a PDF-page inclusion, a table or a TikZ picture; no image file is copied into this package, the package declares no asset dependency and no image file is redistributed with it. Licence: version arXiv:2608.04650v1 of this article is distributed under the Creative Commons Attribution 4.0 International licence, as recorded in the arXiv licence metadata for this exact version and shown on the abstract page https://arxiv.org/abs/2608.04650v1. A full rights notice is delivered at licenses/arxiv-2608.04650-CC-BY-4.0.txt. Characters: every retained excerpt is pure ASCII, so the delivered engine, XeLaTeX with its default Latin Modern text font, reports no missing character.
\begin{lemma}\label{lem:APinfty} Let $G = \Gpres{a, t }{a(tat^{-1})a^{-1}(ta^{-1}t^{-1})= 1 }$. Then the subgroup of $G$ generated by $X_\infty = \{ t^{i} a t^{-i} : i \in \mathbb{Z} \} $ is isomorphic to 
the right-angled Artin group
$A(P_\infty)$ where $P_\infty$ is the graph with vertex set $\Z$ and $i$ adjacent to $i+1$ for all $i \in \mathbb{Z}$. 
Furthermore, for any subset $Y$ of $\Z$ the subgroup of $G$ generated by $X_Y = \{ t^i a t^{-i} : i \in Y \}$ is isomorphic to $A(\Delta)$ where $\Delta$ is the induced subgraph of $P_\infty$ on the set of vertices $Y \subseteq \Z$.     
\end{lemma}

\begin{lemma}\label{lem:UsefulTExpSum} 
Let $G = \Gpres{a, t }{a(tat^{-1})a^{-1}(ta^{-1}t^{-1})= 1 }$ and let 
$U$ be a finite set of words over $\{a,t\}^{\pm 1}$, let $w$ be a   
word over $\{a,t\}^{\pm 1}$ and suppose that $w \in \Mgen{U}$.   
If every word in $U$ has non-negative $t$-exponent sum and if $w$ has $t$-exponent sum zero then $w \in \Mgen{U_0}$ where $U_0 \subseteq U$ is the subset of words with $t$-exponent sum zero.          
\end{lemma}

\begin{theorem}\label{thm:PrefixTech} 
Let $G = \Gpres{a, t }{a(tat^{-1})a^{-1}(ta^{-1}t^{-1})= 1 }$   and set $K = \Mgen{X}$ where $X = \{ t^5at^{-5}, t^{6}at^{-6}, t^7at^{-7}, t^8at^{-8} \}^{\pm 1}$. Let $W = \{w_i: i=1,\ldots, n \}$ be a finite subset of    $X^*$,  and let $\mathcal{P}_W$ be the one-relator group presentation 
\[
\Gpres{a,t}{
a(tat^{-1})a^{-1}(ta^{-1}t^{-1})
\left(\prod_{1 \leq i \leq n}     
(t^3at^{-3})w_i(t^3a^{-1}t^{-3})(t^3at^{-3})w_i^{-1}(t^3a^{-1}t^{-3})
\right)= 1}.
\]
If membership in $\Mgen{W} \leq K$ is undecidable then the one-relator group presentation $\mathcal{P}_W$ has undecidable prefix membership problem.
  \end{theorem}

\begin{lemma}\label{la2}
With the notation from the statement of  Theorem~\ref{thm:PrefixTech}, 
for any $\alpha \in X^{*}$ we have 
\[ 
(t^3at^{-3})\alpha (t^3a^{-1}t^{-3}) \in P_{W} \Longleftrightarrow \alpha \in \Mgen{W}.
\]
\end{lemma}

\begin{theorem}\label{thm:PrefixTech:Reduced} 
Let $G = \Gpres{a, t }{a(tat^{-1})a^{-1}(ta^{-1}t^{-1})= 1 }$   and set $K = \Mgen{X}$ where $X = \{ t^5at^{-5}, t^{6}at^{-6}, t^7at^{-7}, t^8at^{-8} \}^{\pm 1}$.  Fix  
\[ r \equiv a(a(tat^{-1})a^{-1}(ta^{-1}t^{-1}))a^{-1} \] 
and for each word $u \equiv x_1 \ldots x_l \in X^*$ define 
\[
\sigma(u) \equiv rx_1rx_2 \ldots rx_l  r.
\]
Let $W = \{w_i: i=1,\ldots, n \}$ be a finite subset of  $X^*$. For each $i = 1 \ldots n$ set 
\[
\alpha_i \equiv 
(t^3 a t^{-3}) \sigma(w_i) 
(t^3 a^{-1} t^{-3}) r  
(t^3 a t^{-3}) \sigma(w_i^{-1}) 
(t^3 a^{-1} t^{-3}). 
\]
Set $\alpha \equiv \alpha_1 r \alpha_2 r \ldots r \alpha_{n-1} r \alpha_n$. 
Let  $\mathcal{Q}_W$ be the one-relator group presentation 
\[
\Gpres{a,t}{
\alpha r \alpha^{-1} = 1}.
\]
Then $\alpha r \alpha^{-1}$ is a reduced word. Furthermore, if membership in $\Mgen{W} \leq K$ is undecidable then the one-relator group presentation $\mathcal{Q}_W$ has undecidable prefix membership problem.
  \end{theorem}

\begin{lemma}\label{lem:PrefGenSet:Reduced} 
With the notation from the statement of Theorem~\ref{thm:PrefixTech:Reduced}, the prefix monoid $Q_W := \mathrm{pref}(\mathcal{Q}_W)$ is generated by the finite set 
\begin{align*}
\{ a, a^{-1}, \; t, \; tat^{-1}, \; t^3at^{-3} \} \cup 
\{ \mathrm{pref} ((t^3at^{-3})w_i(t^3a^{-1}t^{-3})) : i = 1, \dots, n\}.
\end{align*}
\end{lemma}

\begin{lemma}\label{la3}
With the  notation from the statement of Theorem~\ref{thm:PrefixTech:Reduced}, for any $\alpha \in X^{*}$, we have 
\[ 
(t^3at^{-3})\alpha (t^3a^{-1}t^{-3}) \in Q_{W} \Longleftrightarrow \alpha \in \Mgen{W}.
\]
\end{lemma}

\begin{proof}
($\Longleftarrow$)  Let $\alpha \in \mathrm{Mon}\langle W \rangle$. By Lemma~\ref{la2} we have   
\[
(t^3at^{-3})\alpha (t^3a^{-1}t^{-3}) \in  P_W \subseteq Q_{W},  
\]
as required. 
    
($\Longrightarrow$)  Suppose that $(t^3at^{-3})\alpha (t^3a^{-1}t^{-3}) \in Q_{W}$, where $\alpha \in X^*$.  Observe that the word $(t^3at^{-3})\alpha (t^3a^{-1}t^{-3})$ has $t$-exponent sum zero. 

Since  by Lemma~\ref{lem:PrefGenSet:Reduced} the element $(t^3at^{-3})\alpha (t^3a^{-1}t^{-3})$ can be written as a product of the generators in the statement of that lemma, and since those elements all have non-negative $t$-exponent sum, it follows from  Lemma~\ref{lem:UsefulTExpSum}  that in $G$ the word $(t^3at^{-3})\alpha (t^3a^{-1}t^{-3})$ is equal to a product of words from the set   
\begin{align*}
        & Y : = \{ a, a^{-1}, tat^{-1}, t^3at^{-3} \} \cup \\ 
& \{ \text{prefix}(t^3at^{-3})w_i(t^3a^{-1}t^{-3}) \; \mbox{with $t$-exp sum $0$}, \; i=1, \ldots, n \}. 
    \end{align*}
Observe that any product of words from the set $Y$ must, by definition of $W$, be equal to a product of words from the set    
\[
Z =  \{a, tat^{-1}, t^3at^{-3}, t^5at^{-5}, \dots, t^8at^{-8}\}^{\pm 1}.
\]
Now consider the subgroup $L$ of $G$ generated by this set $Z$. It follows from Lemma~\ref{lem:APinfty} that  
\[
L \cong (\Z \times \Z) \ast \Z \ast A(P_4).
\]
Now arguing in exactly the same way as in the proof of Lemma~\ref{la2}, and using the same notation as in that proof, it follows that   $(t^3at^{-3})\alpha (t^3a^{-1}t^{-3})$ is equal in $G$ to a product of words from $Y_0$ that when written as words over the generators $\{ x_0, y_5, y_6, y_7, y_8 \}^{\pm 1}$ have $x_0$-exponent sum zero. Then by inspection of the words in the set $Y_0$ it follows that $(t^3at^{-3})\alpha (t^3a^{-1}t^{-3})$ is equal in $G$ to a product of words from the set    
\[
\{ (t^3at^{-3})w_i(t^3a^{-1}t^{-3}) \; i=1, \ldots, n \}.
\]
But then, again arguing as in the proof of Lemma~\ref{la2}, it follows from this that  
$\alpha \in \Mgen{W}$, as required.  
\end{proof}

\end{document}
